% Figures 1 and 2 with their claims: text for the paper (#149).
% Every interval is from artifacts/issue-149/figure_claims.md (generated by
% scripts/figure_claims.py with the #140 settings: 95% percentile bootstrap of
% test sources within pairs, paired randomisation p, 10,000 replicates, seed
% 42); D3's overall C3/C10 test and its Holm p-value are from report.md.
% Figures: scripts/generate_figure1_d3_scatter.py and
% scripts/generate_figure2_model_comparison.py (default --metric mrr).

\paragraph{V6 versus V8 on the Conference pairs.}
On D3, adding the key--value description (V2+V6, C3) ranks better than adding
the relational signature (V2+V8, C10) under the same model: MRR
$-0.078$ for C10 [$-0.114$, $-0.042$], Holm-adjusted $p=0.0006$. The gap is
not uniform across pairs (Figure~\ref{fig:d3_scatter}). On the ten pairs that
involve neither CMT nor EDAS, C3 leads by $0.131$ [$0.074$, $0.191$]; on the
eleven CMT or EDAS pairs the two are indistinguishable ($-0.030$
[$-0.075$, $0.016$]). The difference between the two groups, $0.101$, has an
interval of [$0.028$, $0.176$] when test sources are resampled within pairs
and [$0.023$, $0.182$] when whole pairs are resampled within each group.
C10 is ahead on three pairs (cmt-edas, edas-ekaw, edas-sigkdd), but no
pair's own interval favours it, while four favour C3 (conference-sigkdd,
confof-ekaw, ekaw-sigkdd and confof-edas). Hits@1 shows the same group
difference ($0.143$ [$0.036$, $0.251$]), although under Hits@1 C10 is ahead
on four pairs (edas-iasted as well), again none with an interval that
favours it. F1 is weaker. Its group difference ($0.097$ [$-0.020$, $0.212$])
includes zero. For its overall difference ($-0.054$) the bootstrap interval
[$-0.111$, $0.005$] just includes zero while the randomisation test gives
$p=0.049$; the two disagree at the margin, and we draw no conclusion from it.

We do not attribute the group difference to annotation sparsity. The
CMT/EDAS grouping comes from an earlier hypothesis that V8 compensates for
sparsely annotated ontologies, but by annotation density (the mean number
of seven annotation predicates present per entity) EKAW (0.019), not CMT
(0.064) or EDAS (0.083), is the sparsest of the seven ontologies; IASTED,
SIGKDD and Conference lie between 0.090 and 0.114, and only ConfOf (0.627)
is markedly denser. The grouping was fixed before these results, yet the
comparison remains exploratory: each ontology appears in six pairs, so the
pairs are not independent and the pair-level interval is optimistic, and
each pair has only 3--20 test reference sources.

\begin{figure}[t]
  \centering
  \includegraphics[width=0.62\linewidth]{figures/fig1_d3_scatter.pdf}
  \caption{Re-ranking MRR of C3 (V2+V6, M2) against C10 (V2+V8, M2) on the
  21 D3 pairs. Points below the diagonal are pairs where C3 is ahead; orange
  points are pairs involving CMT or EDAS. Labelled: the three pairs where C10
  is ahead, none with an interval that excludes zero.}
  \label{fig:d3_scatter}
\end{figure}

\paragraph{Model choice with the verbalisation fixed.}
Figure~\ref{fig:model_comparison} holds the strategy at V2+V8 and varies the
model on D1 and D2 (all pairwise differences in MRR; Holm within each
dataset's figure comparisons, which are exploratory and separate from the
registered ones). The domain-specific model does not lead. On D1, M5 (a
general-purpose model) is above M4 (BioLORD) by $0.026$ [$0.018$, $0.034$]
and above M3 by $0.019$ [$0.012$, $0.025$]; M3 and M4 do not differ after
adjustment ($-0.0075$ [$-0.0148$, $-0.0003$] unadjusted, Holm $p=0.079$). On D2, M4 and M5
are tied ($0.000$ [$-0.012$, $0.011$]), and both are above M3 (M4 by $0.047$
[$0.033$, $0.060$], M5 by $0.047$ [$0.036$, $0.057$]). M1 and M2 are
indistinguishable on both datasets and below every other model.

No V2+V8 model reaches V1 on M4 (C17). C17 is above the best V2+V8 model on
D1 by $0.045$ [$0.036$, $0.054$] (over M5) and on D2 by $0.089$
[$0.071$, $0.107$] (over M4); with M4 held fixed, V1 is above V2+V8 by
$0.071$ [$0.061$, $0.080$] on D1 and $0.089$ [$0.071$, $0.107$] on D2. On
these two datasets, then, the choice of text moves MRR more than the choice
among these models. V1 differs from V2+V8 by omitting the synonyms, the
definitions and the relational signature, so this does not isolate V8.
Hits@1 agrees. On D1 its interval places M3 below M4
($-0.015$ [$-0.026$, $-0.004$]); MRR points the same way before adjustment
($-0.0075$ [$-0.0148$, $-0.0003$]) and only its Holm-adjusted test leaves the
two unseparated. F1 separates two pairs that MRR ties: M2
above M1 on D1 ($0.024$ [$0.016$, $0.032$]) and M5 above M4 on D2 ($0.044$
[$0.025$, $0.063$]).

\begin{figure}[t]
  \centering
  \includegraphics[width=0.9\linewidth]{figures/fig2_model_comparison.pdf}
  \caption{Re-ranking MRR on D1 (light) and D2 (dark) for each model with
  V2+V8 held constant (C18, C10, C15, C11, C13), and V1 on M4 (C17) for
  reference.}
  \label{fig:model_comparison}
\end{figure}
